\documentclass[11pt,a4paper]{article}
\usepackage[margin=27mm]{geometry}
\usepackage{amsmath,amssymb,amsthm,fontspec,unicode-math,booktabs,hyperref}
\setmainfont{STIXGeneral}\setmathfont{STIX Two Math}
\setlength{\emergencystretch}{3em}
\newtheorem{theorem}{Theorem}
\newtheorem{lemma}[theorem]{Lemma}
\newcommand{\gen}[1]{\langle #1\rangle}
\title{A nonfirm obstruction to an indexed HNN construction}
\author{Human Compression project}
\date{11 October 2026}
\begin{document}
\maketitle
\begin{abstract}
We give a rank-six flag-transitive coset geometry of $S_4$ and an
automorphism preserving its entire unlabelled standard-parabolic lattice.
The literal HNN construction of Piedade--Tranchida, arXiv:2505.24662v1,
Definition 3.31, produces three pairwise incident cosets that cannot extend
to a chamber. Thus the unrestricted assertion of Theorem 3.35 and its
geometry conclusion require additional hypotheses. The input is neither
firm nor residually connected. The proof below is elementary and uniform
in the infinite cyclic coordinate; a finite exhaustive calculation is
separate validation.
\end{abstract}

\section{The exact source boundary}
Throughout, incidence means nonempty intersection of left cosets. The
source's Section 2.1, pages 6--8, permits arbitrary indexed subgroups as
the standard maximal parabolics; this name does not require maximality in
the subgroup lattice. It sets $A_J=\bigcap_{j\in J}A_j$ and
$A_\varnothing=A$. A flag-transitive coset incidence system is a geometry,
but firmness and residual connectedness are separate conditions.

Definition 3.30, pages 27--28, allows $C=A_J$ and $D=A_K$ for arbitrary
$J,K\subseteq I$. An admissible isomorphism $\phi:C\to D$ must preserve
the unlabelled lattices of standard parabolics contained in these groups,
including their intersections. In particular, $J=K=\varnothing$ and
$C=D=A$ are allowed. Theorem 3.35, page 30, assumes only flag transitivity
of the input and this admissibility for its initial assertion. Firmness
and residual connectedness occur as additional hypotheses for subsequent
conclusions. No such hidden assumptions are imposed in the definitions.

For $J=K=\varnothing$, Definition 3.31 identifies named indices whenever
$\phi(A_i)=A_j$, taking the equivalence relation generated by these pairs.
For a resulting class $s$ it prescribes
\[
 G_s=\gen{t,A_s},\qquad A_s=\bigcap_{i\in s}A_i,
 \qquad G_t=A,
 \qquad G=\gen{A,t\mid t^{-1}xt=\phi(x)\ (x\in A)}.
\]
Preservation of the unlabelled lattice does not assert that $\phi$ permutes
the individually named $A_i$. This is the distinction used here.

\section{An elementary input geometry}
In $A=S_{\{1,2,3,4\}}$, put
\[
 a=(34),\quad b=(12),\quad r=(13)(24),\quad s=(14)(23),
 \qquad k=ab=rs=(12)(34).
\]
Let
\[
 P=\gen{(13),(24)},\qquad Q=\gen{(14),(23)},\qquad
 D=\gen{a,b,r},
\]
and use the following six indexed subgroups:
\[
 A_1=\gen a,\quad A_2=\gen b,\quad A_3=P,\quad
 A_4=Q,\quad A_5=D,\quad A_6=\gen r.                 \tag{1}
\]
Here $D$ has order eight and preserves the unordered partition
$\{\{1,2\},\{3,4\}\}$. Its normal forms are $a^i b^j r^\ell$, with
$i,j,\ell\in\{0,1\}$; conjugation by $r$ swaps $a$ and $b$.
The standard parabolics are exactly
\[
 \mathcal L=\{A,D,P,Q,\gen a,\gen b,\gen r,\gen s,1\}.       \tag{2}
\]
Indeed, $D\cap P=\gen r$, $D\cap Q=\gen s$, and $P\cap Q=1$.
All four listed order-two groups lie in $D$; among them only $\gen r$
lies in $P$ and only $\gen s$ lies in $Q$. Distinct order-two groups
intersect trivially. These observations prove intersection closure and
show that every member of (2) is a standard parabolic of (1).

\begin{lemma}\label{lem:helly}
Every pairwise-intersecting family of left cosets of members of
$\mathcal L$ has a common representative.
\end{lemma}
\begin{proof}
Cosets of $A$ impose no restriction. A coset of $1$ is a singleton, so
pairwise intersection with it already gives the conclusion. Repeated
cosets can be removed; distinct cosets of the same subgroup are disjoint.
It remains to consider at most one coset of each of the other seven groups.

Call cosets of $\gen a,\gen b,\gen r,\gen s$ atomic. These are edges
with endpoints $x,xz$, where $z\in\{a,b,r,s\}$. Every atomic coset lies
in one left coset of $D$. Pairwise-intersecting atomic cosets therefore
lie in the same $D$-coset. The homomorphism
\[
 \epsilon:D\longrightarrow C_2,\qquad
 \epsilon(a^i b^j r^\ell)=i+j+\ell\pmod2
\]
takes all four of $a,b,r,s=abr$ to $1$. Thus their Cayley graph on each
$D$-coset is bipartite. Pairwise-intersecting edges in a triangle-free
simple graph have a common vertex: otherwise two edges meeting at $x$
would be met away from $x$ by a third edge, giving a triangle.

Suppose at least two atomic cosets occur, with common vertex $x$.
Their intersection is exactly $\{x\}$. A $D$-coset meeting an atomic
coset contains it, and hence contains $x$. Consider a $P$-coset meeting
all the atomic cosets. If the $\gen r$-coset occurs, it is contained in
that $P$-coset, so again $x$ belongs. Otherwise take two distinct atomic
generators from $a,b,s$. If the $P$-coset omitted $x$, it would contain
both opposite endpoints $xz_1,xz_2$, forcing $z_1z_2\in P$.
The product $ab=k$ is not in $P$, and $as,bs$ have order four, whereas
$P$ has exponent two. This is impossible. For $Q$ the same argument
uses $s$ in place of $r$ and the three remaining generators $a,b,r$.
Thus $x$ lies in every coset in the family.

Suppose exactly one atomic coset occurs. A $D$-coset is then redundant.
For the $\gen r$-coset a $P$-constraint is redundant, and for the
$\gen s$-coset a $Q$-constraint is redundant; the remaining pairwise
constraints have a common point. For an atom of type $a$ or $b$, the only
remaining possible failure would have a $P$-coset meeting one endpoint
and a $Q$-coset meeting the other. Their mutual intersection would imply
$z\in PQ$ (or $QP$) for $z=a$ or $b$. This does not happen. Multiplying
$a$ or $b$ by the four members $1,(14),(23),s$ of $Q$ gives, respectively,
the transposition $a$ or $b$, two three-cycles, and a four-cycle, none in
$P$. Inversion gives the same exclusion for $QP$.

Finally suppose no atomic coset occurs. Only a $P$-, a $Q$-, and a
$D$-coset remain. If all three occur, translate the unique intersection
of the first two to $1$, so those two cosets are $P,Q$. Write the third
as $gD$. It meets both only if $g^{-1}\in DP\cap DQ$. But
\[
 DP=D\cup D(13),\qquad DQ=D\cup D(14),\qquad DP\cap DQ=D.
\]
The two nontrivial right $D$-cosets are distinct, since $(13)(14)$ has
order three and cannot lie in the group $D$ of order eight. Hence
$g\in D$, so the third coset contains $1$ as well.
\end{proof}

Lemma~\ref{lem:helly} proves that every flag for (1) has a common
representative $x$. It is therefore a translate by $x$ of a standard
flag and extends to $\{xA_i:1\leq i\leq6\}$. Thus the input is a
flag-transitive coset geometry, with no reliance on the finite enumeration.
Its Borel is $1$. Every minimal parabolic is $1$, since at least two of
the three distinct groups $\gen a,\gen b,\gen r$ remain when any single
type is deleted. Every standard panel consequently has one chamber,
so the geometry is not firm. It is not residually connected either:
the residue of the standard flag of types $\{1,5\}$ has group $\gen a$
and four remaining subgroups all equal to $1$. Its incidence graph is
two disjoint complete graphs on four vertices.

\section{Admissibility and the literal output}
Define $\phi(x)=bxb$. It fixes $A,D,\gen a,\gen b,1$ and interchanges
$P,Q$ and $\gen r,\gen s$. Hence it preserves exactly the lattice (2)
and all intersections: it satisfies Definition 3.30. The named index
classes in Definition 3.31 are
\[
 \{1\},\quad\{2\},\quad\{3,4\},\quad\{5\},\quad\{6\}.    \tag{3}
\]
In particular, $\phi(A_6)=\gen s$ is a standard parabolic but is not
one of the six named $A_i$.

For clarity, use the following coordinate convention for the infinite
group. Represent $xt^n$ by $(x,n)\in A\times\mathbb Z$ and put
\[
 (x,n)(y,m)=(x\phi^{-n}(y),n+m).                         \tag{4}
\]
Associativity follows from the automorphism law, and
$(1,-1)(x,0)(1,1)=(\phi(x),0)$. Conversely the relation in the HNN
presentation moves every letter from $A$ to the left, giving $xt^n$;
the map into (4) shows uniqueness. Thus (4) is precisely the presented
HNN extension, not merely a quotient of it. In this example
$\phi^{-n}=\phi^n$ because $\phi^2=1$.

Let $V=\gen{r,s}=\{1,r,s,k\}$. For each $\phi$-invariant subgroup
$B\leq A$, the subgroup $\gen{B,t}$ is the set $B\times\mathbb Z$
under (4). The literal maximal parabolics of the output are therefore
\[
\begin{array}{c|c}
\text{type}&\text{subgroup in coordinates (4)}\\\hline
\{1\}&\gen a\times\mathbb Z\\
\{2\}&\gen b\times\mathbb Z\\
\{3,4\}&\{1\}\times\mathbb Z\\
\{5\}&D\times\mathbb Z\\
\{6\}&V\times\mathbb Z\\
t&A\times\{0\}
\end{array}                                                     \tag{5}
\]
Here the sixth input group expands because conjugation by $t$ changes
$r$ to $s$, and $P\cap Q=1$ accounts for the third row.

\begin{theorem}
The output of Definition 3.31 for (1) and $\phi(x)=bxb$ is neither
flag-transitive nor a geometry.
\end{theorem}
\begin{proof}
Write $G_1,G_2,G_6$ for the output groups at the corresponding singleton
classes. The three distinct-type cosets
\[
 G_1,\qquad G_2,\qquad aG_6                              \tag{6}
\]
meet pairwise: their respective intersections contain $1,a,b$, since
$ab=k\in V$. On the other hand,
\[
 G_1\cap G_2=\{1\}\times\mathbb Z=\gen t,
 \qquad aG_6=aV\times\mathbb Z.
\]
The coset $aV$ omits $1$ (its elements are odd permutations), so (6)
has empty total intersection. A translate of the standard flag has a
common representative. Hence this flag cannot be in its orbit, proving
failure of flag transitivity.

More strongly, the output type $\{3,4\}$ has subgroup $\gen t$.
Any one of its cosets incident with both $G_1$ and $G_2$ is contained
in both: $\gen t$ is a subgroup of each. Such a coset must be exactly
$\gen t$, which does not meet $aG_6$. Thus (6) cannot extend to a
chamber. Since the type set is finite, it extends to a maximal flag
that is not a chamber. The output is not a geometry.
\end{proof}

\section{Scope, attribution, and reproducibility}
This is a correction to the unrestricted Theorem 3.35 and clause (a) of
the specified first arXiv edition. The construction and historical
assertions belong to C.~A.~Piedade and P.~Tranchida, \emph{From Group
Operations to Geometric Structures: Amalgamations, HNN-Extensions, and
Twisting in Coset Geometries}, arXiv:2505.24662v1, Definitions 3.30--3.31
and Theorem 3.35, pages 27--31. The example does not refute a theorem
restricted to firm inputs, or to residually connected inputs. It is
also not a refutation of the abstract HNN construction. No conclusion
about a later edition, historical priority, or a numbered open question
is asserted.

The separate program \texttt{literal\_s4.py} reconstructs all permutations
and literal flags without a named-group library. It checks the nine
parabolics, admissibility, classes (3), and the finite fibers of (6).
The flag counts by size are
\[
 (1,51,288,468,360,144,24).
\]
These computations validate the example; Lemma~\ref{lem:helly}, the
coordinate proof (4), and the three explicit intersections supply the
written proof. No Lean formalization is claimed. Full finite-search
traces and source text are not part of the proof package.
\end{document}
